Se observa que los autómatas celulares neuronales han evolucionado desde sus primeras aplicaciones en morfogénesis, donde se exploró su capacidad de generar estructuras complejas partiendo de una única célula activa, hacia modelos que buscan una mayor versatilidad generativa, capaces de producir dinámicas visuales, texturas, comportamientos adaptativos e incluso manipular estrategias, muchos de ellos inspirados en la definición de Mordvintsev. Además, se han desarrollado arquitecturas híbridas que no generan un autómata celular neuronal individual, sino que modelan un espacio de autómatas semejantes, así como configuraciones generativas y bioinspiradas.

En resumen, estos estudios muestran que los autómatas celulares neuronales no solo son un concepto interesante desde el punto de vista teórico, sino que también son útiles en inteligencia artificial, ciencias naturales e ingeniería. Esta diversidad de enfoques abre la puerta a explorar nuevas adaptaciones y propuestas, especialmente en áreas de aprendizaje profundo, autoorganización y simulación de procesos dinámicos.

